\documentclass{article}
\usepackage{amsmath,amssymb}
\usepackage{booktabs}
\usepackage{geometry}
\geometry{margin=1in}

\title{Step 225: CTMT Recursion as a Candidate Foundational Typed Condition}
\author{Six Birds Foundations III}
\date{2026-05-16}

\begin{document}
\maketitle

\section{Status}

This document formalizes CTMT recursion as a candidate foundational typed condition. Its status is
\[
\texttt{candidate (verified-on-5-track-instances) corpus-pending}.
\]
It is not claimed as a formal category-theoretic theorem.

\section{Theorem Statement}

\textbf{Theorem (CTMT Recursion, candidate).}
Let \(C\) be a typed cascade with a CTMT-stuck verdict of the form ``blocked at external classical theorem \(X\)'', produced by a codex-style training-memory-bounded audit.

\begin{enumerate}
\item A manager-led paper-grounded literature audit may resolve \(X\), by identifying a published theorem, formula, construction, estimate, or atlas condition supplying the named content.
\item When \(X\) is inherited, the cascade-internal computation blocked by \(X\) becomes executable.
\item Executing the reduction typically exposes a deeper typed gate \(X_{\rm next}\), rather than closing the parent residual.
\item The CTMT classification persists. Only the specific stuck-at object refines.
\item The recursion continues until reaching a precision-limited terminus, a theorem-unstated-in-literature terminus, a known foundational barrier, or a genuine closure.
\item The terminal object adapts to the mathematical track: matrix elements, component maps, cycle columns, PDE estimates, or package atlases.
\end{enumerate}

\section{Proof Outline}

At cascade level \(n\), the stuck theorem \(X_n\) is a specific external classical statement. A literature audit can promote \(X_n\) from unavailable content to inherited content. Once inherited, the cascade can perform the blocked reduction. But this operation often reveals that \(X_n\) has extra hypotheses, is adjacent rather than exact, or exposes a finer terminal object. Thus \(X_n\) refines to \(X_{n+1}\). The residual frontier moves downward while the typed-condition shape persists.

The process terminates only when a later gate is missing from the literature, numerical precision becomes limiting, a known barrier is reached, or closure genuinely occurs.

\section{Instance Table}

\begin{tabular}{p{0.13\linewidth}p{0.16\linewidth}p{0.11\linewidth}p{0.20\linewidth}p{0.30\linewidth}}
\toprule
Track & Domain & Length & Terminal object & Chain summary\\
\midrule
RH Branch B & operator analytic & 8 & matrix element & \(\kappa\) to diagonal Gram, L'Hopital, HS norm, transport sampling, candidate disagreement, \(G^{-1}\) conditioning\\
RH Branch A & operator analytic & 4 & Calkin / matrix element & \(\kappa\), Calkin symbol, Connes--Consani insufficient, wavepacket Weyl certificate\\
BSD & motivic / p-adic & 6 & component map & bridge defect, Bloch--Kato line, ETNC, Iwasawa, special cases, component obligations\\
Hodge & algebraic geometric & 7 & cycle column & standard residual, repair column, candidate equations, cycle realization, cycle class/projector, HR projection, descent\\
Navier--Stokes & PDE analytic & 10 & PDE estimate & physical-native gap, packing/amplitude, mismatch, derivative stack, Gevrey, radius-window and time gate\\
P-vs-NP & discrete complexity & 11 & package atlas & package residual, fixed quotient, local-status Tseitin, proof-system bridges, atlas gap, no-smuggling and barriers\\
\bottomrule
\end{tabular}

\section{Terminal-Object Adaptation}

The same typed pattern is preserved while the terminal object changes:
\[
\text{external theorem}
\to
\text{literature resolution}
\to
\text{deeper gate}
\to
\text{refined residual frontier}.
\]
The terminal object is matrix-element-terminal on RH, component-map-terminal on BSD, column-terminal on Hodge, estimate-terminal on Navier--Stokes, and package-atlas-terminal on P-vs-NP.

\section{Corpus Recommendation}

The condition should be recorded alongside the Cascade Reduction Theorem, CTMT, CRCFT, the Framework RH-Carrier Dichotomy, and the Bridge Impossibility Corollary. Recommended corpus locations are \texttt{adequacy.tex}, \texttt{needles.tex}, \texttt{paper/sections/foundational\_typed\_conditions.tex}, and \texttt{anti\_loc/findings\_framework.md}.

\section{Verdict}

\[
\boxed{\texttt{V\_CTMT\_recursion\_formalized\_5\_track}}
\]

\end{document}
